% Generated draft: requires mathematical and source review.
\documentclass[11pt]{article}
\usepackage[margin=25mm]{geometry}
\usepackage{fontspec}
\setmainfont{Latin Modern Roman}
\usepackage{amsmath,amssymb,mathtools,unicode-math}
\setmathfont{Latin Modern Math}
\usepackage{enumitem,needspace,xurl,hyperref}
\setlength{\parindent}{0pt}
\setlength{\parskip}{5pt}
\setlength{\emergencystretch}{3em}
\newcommand{\N}{\mathbb{N}}
\newcommand{\Z}{\mathbb{Z}}
\newcommand{\Q}{\mathbb{Q}}
\newcommand{\R}{\mathbb{R}}
\newcommand{\C}{\mathbb{C}}
\newcommand{\F}{\mathbb{F}}
\newcommand{\A}{\mathbb{A}}
\newcommand{\PP}{\mathbb{P}}
\newcommand{\E}{\mathbb{E}}
\newcommand{\eps}{\varepsilon}
\newcommand{\dd}{\,\mathrm{d}}
\begin{document}
\section*{UnsolvedMath ID: 11000059}
\textbf{Source problem code:} catalogue entry
\subsection{Definitions and mathematical statement}
Let $\mathcal I_g$ be the Torelli group of a closed orientable genus-$g$ surface. Write
\[
 H_*(\mathcal I_g;\Q)=\bigoplus_{j\ge0}H_j(\mathcal I_g;\Q)
\]
for its total rational group homology.

\textbf{Problem 5.14.} Extend the result attributed to Akita in the surrounding source to every genus satisfying $2<g<7$, that is, $g=3,4,5,6$.

The extracted record omits the statement of Akita's theorem. The immediately surrounding local records identify it as Theorem 5.13, place it in a discussion of homological infiniteness and the absence of finite classifying complexes for Torelli and Johnson-kernel groups, and follow this problem with an explicit-cycle question in rational homology. Under the natural total-rational-homology interpretation, the requested small-genus target is
\[
 \dim_{\Q}H_*(\mathcal I_g;\Q)=\infty
 \qquad\text{for each }g\in\{3,4,5,6\}.
\]
This display is a contextual reconstruction requiring the omitted theorem for confirmation, not a certified quotation of it. In particular, the exact homology/cohomology and coefficient conventions, and whether any additional assertion for the Johnson kernel $\mathcal K_g$ is part of the theorem to be extended, are not supplied in this record. Infinite-dimensional total homology does not mean that every degree is infinite-dimensional, and is stronger than merely ruling out a finite classifying complex.

The source remarks that Akita's proof produces no explicit homology classes and uses that as motivation for the following separate problem. An explicit-cycle construction is not silently added as a requirement of the extension question itself.
\subsection{Short English statement}
Extend Akita's homological-infiniteness result to genera three through six. The likely target is infinite-dimensional total rational homology of the Torelli group, but the exact omitted theorem and any Johnson-kernel component must be checked.
\subsection{Sources}
\begin{itemize}
\item Local UnsolvedMath registry, permanent ID 11000059, source code catalogue entry; no individual proposer or proposal year supplied.
\item Farb, edited volume \emph{Problems on Mapping Class Groups and Related Topics}, Problem 5.14, PDF page 48; \url{https://www.math.uchicago.edu/~farb/papers/mcgbook.pdf}. The registry records PDF-font-structure extraction and access on 2026-07-29; the volume inventory distinguishes 329 labeled headings by source page.
\item Adjacent local IDs 11000057, 11000058, and 11000060 refer respectively to Akita's Theorem 5.13, the nonexistence of finite classifying complexes in this context, and explicit independent rational homology classes for $\mathcal I_g$ and $\mathcal K_g$. They do not reproduce Theorem 5.13 itself.
\item Akita is named in the source and summary, but no complete primary bibliographic citation is supplied for the result being extended.
\end{itemize}
\end{document}
